\section{Introduction}
\label{sec:introduction}

Research papers often tell their readers where a study falls short and what should be investigated
next. In clinical research, self-acknowledged limitations help readers interpret findings in light
of study design, data, and applicability~\citep{kilicoglu2018automatic}. In other disciplines,
future-work statements identify extensions, unresolved questions, and further evaluations that
authors consider worth pursuing~\citep{zhu2019recognizing,zhang2023fws}. We refer to these passages
as \emph{author-stated research gaps}: limitations and future-work directions expressed about the
authors' own work. Their extraction provides a way to make this information searchable without
requiring a model to invent a criticism or a research direction.

A collection of such statements can support several activities. A researcher may look for
recurring restrictions in evaluation, identify settings that authors have left untested, or
compare how studies describe future directions. A reviewer may use extracted statements to
locate the authors' discussion of evidence and applicability. Corpus-level work already analyzes
limitation reporting and future-work categories, and traces whether later publications engage
with earlier research directions~\citep{lan2024automatic,hao2020acl,suray2024future}. These uses
motivate a reliable extraction layer, but their interpretation remains conditional: an author's
statement records a reported limitation or direction, not an independently established open
problem. Later work may have addressed it, and the author may not have acknowledged every
important weakness.

The first challenge is therefore definitional. A sentence criticizing prior research can be
useful in a literature review yet lie outside the extraction of a study's own limitations.
Similarly, a statement about how the current paper extends earlier work need not describe
future research. Annotation schemes distinguish these functions in different ways. Argumentative
zoning and scientific-concept annotation include categories related to gaps, failures, and future
work~\citep{liakata2010corpora}, while full-text resources annotate limitations and rhetorical
roles in separate layers~\citep{fisas2016multilayered}. Sci-Sentence provides explicit
Research Gap, Limitation, and Extension categories for literature-review analysis~\citep{bolanos2025scisentence}.
Their names are related to our objective, but their definitions and attribution targets are not
interchangeable. An extraction study must state which source claims count before comparing methods.

The second challenge is structural. Authors may place limitations in a dedicated section,
within a discussion or conclusion, or alongside a particular methodological choice. Dedicated
sections facilitate localization, but their surrounding text also contains results, mitigation,
contribution claims, and descriptions of other studies. Restricting attention to terminal
sections can remove irrelevant material cheaply, while also excluding mid-paper statements.
PDF reading order, broken hyphenation, and sentence boundaries add a further source of error.
A correct sentence classifier cannot recover a gold statement that the parser or localizer has
already omitted. Evaluating localization and detection separately is therefore necessary to
understand the recall of the complete system.

Prior extraction work provides useful precedents. Clinical limitation recognition uses headings,
lexical cues, and supervised or semi-supervised classifiers~\citep{kilicoglu2018automatic}.
Later work on randomized controlled trials (RCTs) separates limitation-sentence detection from
classification into limitation types~\citep{lan2024automatic}. Future-research recognition in
JASIST articles and future-work recognition in NLP similarly separate locating relevant sentences
from grouping or classifying their content~\citep{zhu2019recognizing,zhang2023fws}.
These studies motivate the use of inexpensive structural and lexical evidence before more
flexible semantic classification. They do not imply that one detector transfers equally well
between biomedical full articles, NLP conclusions, and literature-review sentences.

Large language models (LLMs) offer another route, but they introduce a distinct allocation problem.
A full-article prompt presents many sentences unrelated to the extraction target, whereas a
candidate-only prompt depends on the recall of the earlier stages. LLM-based work also combines
source extraction with generation: LimGen uses author-written limitation sections as reference
material, BAGELS evaluates limitation extraction and generation, and FutureGen extracts and
then generates future-work text~\citep{faizullah2024limgen,alazher2025bagels,alazher2025futuregen}.
Critical-limitation and unstated-limitation studies instead ask models to identify weaknesses
that may not have been explicitly acknowledged~\citep{xu2025llms,mishra2026tree}.
The ability to produce a plausible criticism does not establish the ability to identify an
existing author-stated sentence. Our objective is the latter, with source attribution retained.

We investigate a \emph{selective verification funnel} that divides this extraction task into
three stages. Stage~A selects candidate regions using document structure. Stage~B combines sentence
detection, with lexical cues in the hybrid configuration. Stage~C optionally asks an LLM to verify the surviving candidates
against the task's attribution policy. Each stage selects existing parsed sentences rather than
rewriting them. This design makes the consequences of compression visible: we can measure how
many labelled positives remain after localization, how many candidates the detector passes on,
and which true and false positives the verifier removes.

The present experiments examine this mechanism under several distinct input conditions. The
principal analysis uses reconstructed full-article SAL records from biomedical RCT reports, with
complete saved detector predictions and verifier decisions. Separate studies concern curated
future-work sentences in NLP, a literature-gap classification task, and a small arXiv development
reference. Development studies further compare localizer variants, lexical baselines, class
balancing, and alternative reference sets. A documented model-labelled corpus and its implemented
human-validation workflow are distinguished from completed human evaluation. We also compare estimated LLM input and a recorded small full-article prompting baseline.
Dataset-specific labels, prevalence, and processing units are kept explicit because pooling them
would obscure what each experiment measures. These experiments concern extraction and its
processing requirements; corpus-level characterization is a separate empirical objective.

\paragraph{Research questions.}
We ask three questions for the evidence reported in this paper:
\begin{itemize}
  \item \textbf{RQ1 (Stage contribution).} How do structural localization, sentence detection,
        and selective LLM verification change limitation-detection precision, recall, and $F_1$
        on a full-article stream?
  \item \textbf{RQ2 (Detector and reference dependence).} How does the observed trade-off vary
        across detector families, annotation targets, and reference construction, including
        author-stated future work and literature-gap recognition?
  \item \textbf{RQ3 (Input efficiency).} How much LLM input is avoided by selecting candidate
        sentences, and what effectiveness is recorded in a same-model full-article comparison?
\end{itemize}

\paragraph{Contributions.}
The study makes three contributions:
\begin{enumerate}
  \item An extractive pipeline with explicit interfaces for localization, detection, and
        verification, together with the rationale for allocating structural, lexical, learned,
        and LLM processing to different stages.
  \item A full-article stage analysis recomputed from saved sentence-level predictions and
        verifier decisions, accompanied by detector comparisons and separate auxiliary-task
        results that make the input condition and attribution target visible, alongside development
        ablations and reference-sensitivity analysis.
  \item An input-token accounting study and a small same-model full-article comparison,
        reporting effectiveness and processing-volume trade-offs separately from unmeasured
        total billing or throughput.
\end{enumerate}
Prompt development used benchmark feedback, so the verifier results are exploratory.
The claim is a measured trade-off for the examined configurations, rather than general superiority
of LLMs, independent validation across all disciplines, or demonstrated idea-generation quality.

\paragraph{Paper structure.}
\Cref{S:related-works} develops the literature comparison, including the chronological
study-by-study table. \Cref{S:methodology} traces data acquisition, preparation, extraction,
scoring, and the analyses used to answer each RQ. \Cref{S:results} reports the verified and
aggregate results with their evidence status. \Cref{S:discussion} interprets the mechanisms and
validity limits, and \Cref{S:conclusion} summarizes the answers and the next evaluation steps.
